\documentclass[10pt,a4paper]{amsart}
\usepackage[margin=19mm]{geometry}
\usepackage{amsmath,amssymb,mathtools}
\usepackage[scr=rsfs,cal=euler]{mathalfa}
\usepackage{xfrac}
\usepackage[unicode,hidelinks,bookmarks=false]{hyperref}
\newcommand{\ie}{i.e.\ }
\newcommand{\cf}{cf.\ }
\newcommand{\resp}{respectively\ }
\newcommand{\Subsection}{\subsection}
\providecommand{\nobreakdash}{\nobreak\hbox{-}}
\setlength{\emergencystretch}{2em}
\multlinegap0pt
\usepackage{amssymb}
\usepackage{tikz}
\usepackage[shortlabels]{enumitem}
\usepackage{xcolor}
\newtheorem{lemma}{Lemma}
\newtheorem{theorem}{Theorem}
\title{Cyclotomic skew partial difference sets and partial difference families}
\author{Sophie Huczynska, Tekg\"ul Kalayc{\i}}
\date{Source: arXiv:2605.19596v1 (CC BY 4.0)}
\begin{document}
\maketitle

\noindent\emph{Source and licence.} Cyclotomic skew partial difference sets and partial difference families. Version arXiv:2605.19596v1 (CC BY 4.0), distributed by arXiv under the Creative Commons Attribution 4.0 International licence, http://creativecommons.org/licenses/by/4.0/, as recorded in the arXiv licence metadata for this exact version and shown on the abstract page https://arxiv.org/abs/2605.19596v1. Full licence notice: \texttt{licenses/arxiv-2605.19596-CC-BY-4.0.txt}.\par\smallskip

\noindent\textit{Editorial scope.} Counted unit: the article's theorem thm:EPDF1111 (source0.tex lines 1070--1088). The article's complete original proof of it is delivered with it (source0.tex lines 1090--1216, one \texttt{proof} environment of 127 lines). No mathematics is written, repaired or re-derived: the statement and the proof are the article's own lines, in the article's own notation, and no retained line is substituted or re-flowed.  Retained uncounted as the source-local context the proof needs: lines 202--208; lines 1593--1684.  Nothing was omitted from any retained span, so the package carries no omission record.  No retained line contains a citation, so the wrapper carries no bibliography and no bibliography entry.  No retained line contains a figure environment, an image-inclusion command or drawing code, so no image is delivered and the package declares no asset dependency.  Editorial formatting only, in addition: an AMS \texttt{amsart} preamble that restates the article's own declarations and macros used by the retained lines, namely the environments \texttt{lemma}, \texttt{theorem} on separate counters, together with the standard packages those lines need.  The retained environments print in this document as Lemma 1, Theorem 1 in order of appearance, whereas the article prints the same items with the numbering of its own full document.  The complete native source of the article as distributed in the arXiv e-print for this exact version is bundled unchanged as \texttt{sources/200927/source0.tex}.
\begin{lemma}\label{lemma:cyc}
Let $q=ef+1$ be a prime power.  For $0 \leq j,l \leq e-1$:
\begin{itemize}
\item[(i)] $\Delta(C_j^e)=\sum_{i=0}^{e-1} (i,0)_e C_{i+j}^e$
\item[(ii)] $\Delta(C_{j+l}^e,C_l^e)=\sum_{i=0}^{e-1} (i,j)_e C_{i+l}^e$.
\end{itemize}
\end{lemma}

\begin{theorem}\label{thm:e=8cyclo}
Let $q=8f+1$ be a prime power.  The cyclotomic numbers are uniquely determined by $x,y,a$ and $b$, defined as follows.
\begin{itemize}
    \item $q=x^2+4y^2$, $x \equiv 1 \mod 4$ is the unique proper representation of $q=p^m$ if $p \equiv 1 \mod 4$; otherwise $x=\pm p^{m/2}$ and $y=0$.
    \item $q=a^2+2b^2$, $a \equiv 1 \mod 4$, is the unique proper representation of $q=p^m$ if $p \equiv 1$ or $3 \mod 8$; otherwise $a= \pm p^{m/2}$ and $b=0$. 
\end{itemize}
The signs of $y$ and $b$ are ambiguously determined.

\begin{itemize}
\item[]{\bf Case when $q=8f+1$, $f$ odd (i.e $q \equiv 9 \mod 16$)}\\
The relations between the cyclotomic numbers are given by:
\begin{itemize}
\item[]$A=(0,0)_8 =(4,0)_8 =(4,4)_8$; $B=(0,1)_8 =(3,7)_8 =(5,4)_8$;
\item[]$C=(0,2)_8 =(2,6)_8 =(6,4)_8$; $D=(0,3)_8 =(1,5)_8 =(7,4)_8$; 
\item[]$E=(0,4)_8$; $F=(0,5)_8 =(1,4)_8 =(7,3)_8$; 
\item[]$G=(0,6)_8 =(2,4)_8 =(6,2)_8$; $H=(0,7)_8 =(3,4)_8 =(5,1)_8$;
\item[]$I=(1,0)_8 =(3,3)_8 =(4,1)_8 =(4,5)_8 =(5,0)_8 =(7,7)_8$;
\item[]$J=(1,1)_8 =(3,0)_8 =(4,3)_8 =(4,7)_8 =(5,5)_8 =(7,0)_8$;
\item[]$K=(1,2)_8 =(2,7)_8 =(3,6)_8 =(5,3)_8 =(6,5)_8 =(7,1)_8$;
\item[]$L=(1,3)_8 =(1,6)_8 =(2,5)_8 =(6,3)_8 =(7,2)_8 =(7,5)_8$;
\item[]$M=(1,7)_8 =(2,3)_8 =(3,5)_8 =(5,2)_8 =(6,1)_8 =(7,6)_8$;
\item[]$N=(2,0)_8 =(2,2)_8 =(4,2)_8 =(4,6)_8 =(6,0)_8 =(6,6)_8$;
\item[]$O=(2,1)_8 =(3,1)_8 =(3,2)_8 =(5,6)_8 =(5,7)_8 =(6,7)_8$.
\end{itemize}

$
\begin{array}{|c|c|c|}
\hline
q \equiv 9 \mod 16 &\text{$2$ a quartic residue} &\text{$2$ not a quartic residue}\\
\hline
64A & q-15-2x & q-15-10x-8a \\
64B & q+1+2x-4a+16y & q+1+2x-4a-16b\\
64C & q+1+6x+8a-16y& q+1-2x+16y\\
64D & q+1+2x-4a-16y & q+1+2x-4a-16b\\
64E & q+1-18x & q+1+6x+24a\\
64F & q+1+2x-4a+16y & q+1+2x-4a+16b\\
64G & q+1+6x+8a+16y & q+1-2x-16y\\
64H & q+1+2x-4a-16y & q+1+2x-4a+16b\\
64I & q-7+2x+4a & q-7+2x+4a+16y\\
64J & q-7+2x+4a & q-7+2x+4a-16y\\
64K & q+1-6x+4a+16b & q+1+2x-4a\\
64L & q+1+2x-4a & q+1-6x+4a\\
64M & q+1-6x+4a-16b & q+1+2x-4a\\
64N & q-7-2x-8a & q-7+6x\\
64O & q+1+2x-4a & q+1-6x+4a\\
\hline
\end{array}
$

\item[]\noindent{\bf Case when $q=8f+1$, $f$ even (i.e $q \equiv 1 \mod 16$)}\\
The relations between the cyclotomic numbers are given by:
\begin{itemize}
\item[]$A=(0,0)_8$; $B=(0,1)_8 =(1,0)_8 =(7,7)_8$;
\item[]$C=(0,2)_8 =(2,0)_8 =(6,6)_8$; $D=(0,3)_8 =(3,0)_8 =(5,5)_8$; 
\item[]$E=(0,4)_8=(4,0)_8=(4,4)_8$; $F=(0,5)_8 =(5,0)_8 =(3,3)_8$; 
\item[]$G=(0,6)_8 =(6,0)_8 =(2,2)_8$; $H=(0,7)_8 =(7,0)_8 =(1,1)_8$;
\item[]$I=(1,2)_8 =(2,1)_8 =(1,7)_8 =(7,1)_8 =(6,7)_8 =(7,6)_8$;
\item[]$J=(1,3)_8 =(3,1)_8 =(2,7)_8 =(7,2)_8 =(5,6)_8 =(6,5)_8$;
\item[]$K=(1,4)_8 =(4,1)_8 =(3,7)_8 =(7,3)_8 =(4,5)_8 =(5,4)_8$;
\item[]$L=(1,5)_8 =(5,1)_8 =(3,4)_8 =(4,3)_8 =(4,7)_8 =(7,4)_8$;
\item[]$M=(1,6)_8 =(6,1)_8 =(2,3)_8 =(3,2)_8 =(5,7)_8 =(7,5)_8$;
\item[]$N=(2,4)_8 =(4,2)_8 =(2,6)_8 =(6,4)_8 =(4,6)_8 =(6,2)_8$;
\item[]$O=(2,5)_8 =(5,2)_8 =(3,5)_8 =(5,3)_8 =(3,6)_8 =(6,3)_8$.
\end{itemize}

$
\begin{array}{|c|c|c|}
\hline
q \equiv 1 \mod 16 &\text{$2$ a quartic residue} &\text{$2$ not a quartic residue}\\
\hline
64A & q-23-18x-24a & q-23+6x \\
64B & q-7+2x+4a+16y+16b & q-7+2x+4a\\
64C & q-7+6x+16y& q-7-2x-8a-16y\\
64D & q-7+2x+4a-16y+16b & q-7+2x+4a\\
64E & q-7-2x+8a & q-7-10x\\
64F & q-7+2x+4a+16y-16b & q-7+2x+4a\\
64G & q-7+6x-16y & q-7-2x-8a+16y\\
64H & q-7+2x+4a-16y-16b & q-7+2x+4a\\
64I & q+1+2x-4a & q+1-6x+4a\\
64J & q+1-6x+4a & q+1+2x-4a-16b\\
64K & q+1+2x-4a  & q+1+2x-4a+16y\\
64L & q+1+2x-4a & q+1+2x-4a-16y\\
64M & q+1-6x+4a & q+1+2x-4a+16b\\
64N & q+1-2x & q+1+6x+8a\\
64O & q+1+2x-4a & q+1-6x+4a\\
\hline
\end{array}
$

\end{itemize}

\end{theorem}

\begin{theorem}\label{thm:EPDF1111}
	Let $q=8f+1$ be a prime power with $f$ even, and suppose that $q=x^2+4y^2=a^2+2b^2$ are the unique proper representations of $q$. Suppose that $2$ is not a
	quartic residue. Then
	$
	 D=\{C_0^8,C_1^8,C_4^8,C_5^8\}
	$ 
	is:
	\begin{itemize}
		\item a
		$
		(q,4,(q-1)/8;(3q-3+8y)/16,(3q-3-8y)/16)
		$-EPDF relative to $C_0^2$, if $a=-3$ and $y\ne0$;
		
		\item a
		$
		(q,4,(q-1)/8,(3q-3)/16)
		$-EDF, if $a=-3$ and $y=0$.
	\end{itemize}
\end{theorem}

\begin{proof}
	We have
	\[
	\begin{aligned}
		\mathrm{Ext}(\mathcal D)
		={}&
		\Delta(C_0^8,C_1^8)+\Delta(C_1^8,C_0^8)
		+\Delta(C_0^8,C_4^8)+\Delta(C_4^8,C_0^8)\\
		&+\Delta(C_0^8,C_5^8)+\Delta(C_5^8,C_0^8)
		+\Delta(C_1^8,C_4^8)+\Delta(C_4^8,C_1^8)\\
		&+\Delta(C_1^8,C_5^8)+\Delta(C_5^8,C_1^8)
		+\Delta(C_4^8,C_5^8)+\Delta(C_5^8,C_4^8).
	\end{aligned}
	\]
	By Lemma~\ref{lemma:cyc},
	$
	\Delta(C_{j+l}^8,C_l^8)
	=
	\sum_{i=0}^7(i,j)_8C_{i+l}^8.
	$
	Hence each element of $C_i^8$ occurs in $\mathrm{Ext}(\mathcal D)$ precisely
	\[
	\begin{aligned}
		&N_i=(i-1,7)_8+(i,1)_8+(i-4,4)_8
		+(i,4)_8+(i-5,3)_8+(i,5)_8\\
		&\quad +(i-4,5)_8+(i-1,3)_8+(i-5,4)_8
		+(i-1,4)_8+(i-5,7)_8+(i-4,1)_8
	\end{aligned}
	\]
	times.
	
	By Theorem~\ref{thm:e=8cyclo}, we have
	$
	(0,1)_8=(7,7)_8,
	(0,4)_8=(4,4)_8,
	(0,5)_8=(3,3)_8,
	$
	$
	(3,7)_8=(4,1)_8=(4,5)_8=(7,3)_8,
	$
	and
	$
	(3,4)_8=(7,4)_8$.
	So $N_0=N_4=
	2(0,1)_8+2(0,4)_8+2(0,5)_8
	+4(3,7)_8+2(3,4)_8.
	$
	Similarly, using
	$
	(1,3)_8=(6,5)_8,
	(1,7)_8=(2,1)_8,
	$
	$
	(2,4)_8=(5,7)_8=(6,4)_8,
	(2,5)_8=(5,3)_8,
	$
	and
	$
	(1,6)_8=(6,1)_8,
	$
	we have
	$
	N_2=N_6=
	2(1,3)_8+2(3,7)_8+2(1,7)_8
	+2(2,4)_8+2(2,5)_8+2(1,6)_8
	$.
	For the cyclotomic numbers, by Theorem~\ref{thm:e=8cyclo}, since $2$
	is not a quartic residue, we obtain
	$
	N_0=(12q-36-8a+32y)/64
	$
	and
	$
	N_2=(12q+12+8a+32y)/64.
	$
	Thus, if $a=-3$, then
	$
	N_0=N_2=(3q-3+8y)/16.
	$
	
	On the other hand, using
	$
	(0,3)_8=(5,5)_8,
	(0,7)_8=(1,1)_8,
	$
	$
	(3,7)_8=(5,1)_8,
	$
	and
	$
	(3,4)_8=(4,3)_8=(4,7)_8=(1,5)_8$,
	we have
    $N_1=N_5=
	2(0,3)_8+2(0,4)_8+2(0,7)_8
	+2(3,7)_8+4(3,4)_8
	$.
	Using
	$
	(2,3)_8=(3,2)_8=(5,7)_8=(7,5)_8,
	(2,7)_8=(3,1)_8,
	$
	$
	(3,4)_8=(7,4)_8,
	(3,5)_8=(6,3)_8,
	$
	and
	$
	(6,7)_8=(7,1)_8$,
    we have 
	$N_3=N_7=
	2(2,3)_8+2(2,4)_8+2(2,7)_8
	+2(3,4)_8+2(3,5)_8+2(6,7)_8
	$.
	Again by Theorem~\ref{thm:e=8cyclo}, we obtain
	$
	N_1=(12q-36-8a-32y)/64
	$
	and
	$
	N_3=(12q+12+8a-32y)/64.
	$
	Thus, if $a=-3$, then
	$
	N_1=N_3=(3q-3-8y)/16$.
If $y=0$, then the two frequencies coincide and every nonzero element of $\mathbb F_q$ occurs exactly
$(3q-3)/16$ times in $\mathrm{Ext}(\mathcal D)$. 
\end{proof}

\end{document}
